\documentclass[10pt,a4paper]{article}
\usepackage[a4paper,margin=14mm,top=12mm,bottom=14mm]{geometry}
\usepackage{fontspec}
\setmainfont{Liberation Sans}
\usepackage{polyglossia}
\setdefaultlanguage{french}
\usepackage{amsmath}
\usepackage{xcolor}
\usepackage{graphicx}
\usepackage[most]{tcolorbox}
\usepackage{tabularx,booktabs,array,multirow}
\usepackage{enumitem}
\usepackage{fancyhdr}
\usepackage{lastpage}
\usepackage[hidelinks]{hyperref}
\definecolor{uimm}{HTML}{005677}
\definecolor{cyan}{HTML}{13A2DD}
\definecolor{fond}{HTML}{E6F2FA}
\definecolor{ok}{HTML}{2E7D32}
\definecolor{ko}{HTML}{B3261E}
\definecolor{violet}{HTML}{7B1FA2}
\setlength{\parindent}{0pt}
\setlength{\parskip}{3pt}
\setlist{nosep}
\setlist[itemize]{leftmargin=13pt,label=\textcolor{cyan}{\textbullet}}
\newcommand{\titre}[1]{\vspace{5pt}{\color{uimm}\large\bfseries #1}\par\vspace{-3pt}{\color{cyan}\rule{\linewidth}{1.2pt}}\par\vspace{0pt}}
\tcbset{boite/.style={colback=fond,colframe=uimm,boxrule=0.6pt,arc=2pt,left=6pt,right=6pt,top=4pt,bottom=4pt},
        defi/.style={colback=violet!6,colframe=violet,boxrule=0.8pt,arc=2pt,left=6pt,right=6pt,top=4pt,bottom=4pt,
                     title={\bfseries Défi},coltitle=white,colbacktitle=violet,fonttitle=\small}}
\newcommand{\bandeau}[3]{%
\begin{tcolorbox}[colback=uimm,colframe=uimm,arc=2pt,left=4pt,right=8pt,top=4pt,bottom=4pt]
\begin{minipage}[c]{18mm}\setlength{\fboxsep}{1.2pt}\fcolorbox{white}{white}{\includegraphics[height=16.5mm]{logo_uimm.jpg}}\end{minipage}\hspace{4mm}%
\begin{minipage}[c]{\dimexpr\linewidth-22mm\relax}\color{white}
{\small #1}\\[1pt]{\LARGE\bfseries #2}\\[2pt]{\large #3}\hfill{\small CQPM Technicien qualité}
\end{minipage}
\end{tcolorbox}}
\newcommand{\reponse}[1][1]{\par\vspace{2pt}{\color{gray!60}\foreach \i in {1,...,#1}{\rule{\linewidth}{0.3pt}\par\vspace{9pt}}}}
\usepackage{pgffor}
\pagestyle{fancy}\fancyhf{}
\renewcommand{\headrulewidth}{0pt}
\newcommand{\pied}[1]{\fancyfoot[L]{\footnotesize\color{uimm}#1}\fancyfoot[R]{\footnotesize\color{uimm}Page \thepage{} sur \pageref{LastPage}}}

\usepackage{tikz}
\usepackage{longtable}
\renewcommand{\bandeau}[3]{%
\begin{tcolorbox}[colback=uimm,colframe=uimm,arc=2pt,left=4pt,right=8pt,top=4pt,bottom=4pt]
\begin{minipage}[c]{18mm}\setlength{\fboxsep}{1.2pt}\fcolorbox{white}{white}{\includegraphics[height=16.5mm]{logo_uimm.jpg}}\end{minipage}\hspace{4mm}%
\begin{minipage}[c]{\dimexpr\linewidth-22mm\relax}\color{white}
{\small #1}\\[1pt]{\LARGE\bfseries #2}\\[2pt]{\large #3}
\end{minipage}
\end{tcolorbox}}
% en-tête d'une fiche : code, titre, question, bloc, durée, organisation, niveau
\newcommand{\entete}[7]{%
\begin{tcolorbox}[colback=uimm,colframe=uimm,arc=2pt,left=4pt,right=6pt,top=4pt,bottom=4pt]
\begin{minipage}[c]{18mm}\setlength{\fboxsep}{1.2pt}\fcolorbox{white}{white}{\includegraphics[height=16.5mm]{logo_uimm.jpg}}\end{minipage}\hspace{3mm}%
\begin{minipage}[c]{\dimexpr\linewidth-58mm\relax}\color{white}\raggedright
{\small Base SPC · #4}\\[1pt]{\Large\bfseries #1 · #2}\\[2pt]{\itshape #3}
\end{minipage}\hfill
\begin{minipage}[c]{33mm}\color{white}\small\raggedleft
Durée : \textbf{#5}\\ \textbf{#6}\\ Niveau #7 sur 3
\end{minipage}
\end{tcolorbox}}
\newcommand{\tagE}{{\setlength{\fboxsep}{1.5pt}\colorbox{cyan}{\color{white}\scriptsize\bfseries Ellistat}}}
\newcommand{\tagW}{{\setlength{\fboxsep}{1.5pt}\colorbox{uimm}{\color{white}\scriptsize\bfseries Page web}}}
\tcbset{pred/.style={colback=orange!6,colframe=orange!80!black,boxrule=0.6pt,arc=2pt,left=6pt,right=6pt,top=3pt,bottom=1pt,
                    title={\bfseries Je prédis, avant de commencer},coltitle=white,colbacktitle=orange!80!black,fonttitle=\small},
        retenir/.style={colback=fond,colframe=uimm,boxrule=0.6pt,arc=2pt,left=6pt,right=6pt,top=3pt,bottom=3pt,
                    title={\bfseries À retenir},coltitle=white,colbacktitle=uimm,fonttitle=\small},
        plus/.style={colback=violet!6,colframe=violet,boxrule=0.6pt,arc=2pt,left=6pt,right=6pt,top=3pt,bottom=3pt,
                    title={\bfseries Pour aller plus loin},coltitle=white,colbacktitle=violet,fonttitle=\small},
        act/.style={colback=uimm!10,colframe=uimm,boxrule=0.5pt,arc=2pt,left=5pt,right=5pt,top=2pt,bottom=2pt,before skip=7pt,after skip=3pt},
        verif/.style={colback=orange!6,colframe=orange!80!black,boxrule=0.6pt,arc=2pt,left=6pt,right=6pt,top=3pt,bottom=3pt,
                    title={\bfseries À vérifier sur poste avant la première séance},coltitle=white,colbacktitle=orange!80!black,fonttitle=\small}}
\newcommand{\fichier}[1]{\texttt{\small #1}}
\newcommand{\menu}[1]{\textsf{\textbf{#1}}}
\newcommand{\champ}[1]{\rule{#1}{0.3pt}}

\pied{Base SPC · fiche M11 · L'opérateur compte-t-il ?}
\begin{document}
\thispagestyle{fancy}
\entete{M11}{L'opérateur compte-t-il ?}{Trois ou quatre personnes mesurent les mêmes butées : trouvent-elles la même chose ?}{Le moyen de mesure}{70 min}{en groupe}{3}
\begin{tcolorbox}[boite]\textbf{Le but.} Mener une étude R\&R complète : 10 pièces, 3 ou 4 opérateurs, 2 séries.\par\textbf{Il vous faut.} 10 butées numérotées, de diamètres étalés dans la tolérance · Un seul micromètre 0-25 mm · La fiche « L'opérateur compte-t-il ? » (3 pages)\end{tcolorbox}
\begin{tcolorbox}[pred]Qui va peser le plus dans le flou de la mesure : l'instrument ou les opérateurs ?\reponse\end{tcolorbox}
\titre{À faire}
\begin{enumerate}[label=\textcolor{cyan}{\textbf{\arabic*.}},leftmargin=15pt,itemsep=1.5pt]
\item Suivre la fiche « L'opérateur compte-t-il ? » : tirage de l'ordre, série 1, série 2, calculs à la main. La fiche est écrite pour trois opérateurs : à quatre, le quatrième opérateur (D) a ses propres colonnes, comme A, B et C. Chacun note pour le suivant.
\item \tagW\ Dans la page « L'opérateur compte-t-il ? », choisir d'abord le nombre d'opérateurs (3 ou 4) : un changement efface la grille. Saisir les 20 valeurs de chaque opérateur (60 en tout à trois, 80 à quatre) et lire le \%GRR.
\item \tagE\ \menu{Gage R\&R}, onglet \menu{Mesures}, bouton « ajouter un gage R\&R ». Remplir : Nom ; Produits 10 ; Opérateurs 3 ou 4 (le nombre de ceux qui ont mesuré) ; Répétitions 2 (Ellistat propose 3) ; Type de tolérance Bilatéral, Tolérance min 11,85 et max 12,00 ; Méthode de calcul \textbf{ANAVAR} ; K 3 (AIAG). Valider.
\item \tagE\ Cliquer sur « Voir Plus », à côté de « Informations du gage » : vérifier que Tolerance min et Tolerance max sont remplies. Sinon, les saisir (le point sert de virgule).
\item \tagE\ Onglet \menu{Mesure} : un tableau par opérateur, avec les lignes Pièce 1 à Pièce 10 et les colonnes 1 et 2 (les deux répétitions). Saisir les mesures de chaque opérateur ; MOYENNE et ÉTENDUE se calculent seules.
\item \tagE\ Onglet \menu{Analyse} : lire le tableau « Résultats » (répétabilité, reproductibilité, instrument), le bloc « Calcul à partir des tolérances » (GRR\%, Cpc) et, dans « Calcul à partir de la dispersion process », le « Nombre de catégories distinctes ».
\end{enumerate}
\titre{J'observe, je décide}
\begin{enumerate}[label=\textcolor{cyan}{\textbf{Q\arabic*.}},leftmargin=22pt]
\item \%GRR de la page = \champ{14mm} \quad GRR\% d'Ellistat = \champ{14mm} \quad Cpc = \champ{10mm} \quad ndc = \champ{10mm}\par\vspace{5pt}
\item Les deux \%GRR sont-ils égaux ? Sinon, cherchez la raison avec le formateur.\reponse[2]
\item Le système de mesure est-il capable, capable sous conditions ou non capable ? Ellistat et le livret Bosch donnent-ils le même avis ?\reponse[1]
\end{enumerate}
\begin{tcolorbox}[retenir]Le GRR réunit l'instrument (répétabilité) et les opérateurs (reproductibilité). Jusqu'à 10 \% de la tolérance : capable ; au-delà de 30 \% : non capable.\end{tcolorbox}
\begin{tcolorbox}[plus]Dans « Voir Plus », passez K de 3 (AIAG) à 2 (VDA). Que deviennent GRR\% et Cpc ? Que dit ce choix de la façon de compter la dispersion ? Remettez K à 3.\end{tcolorbox}
\vfill{\footnotesize\color{uimm}\textbf{Ressources de la base :} L'opérateur compte-t-il ? (page) · Fiche R\&R (PDF) · Exemple du livret Bosch (3 opérateurs) · Jeu de secours à 4 opérateurs.}
\end{document}